\subsection{Hierarchies from Fluxes in String Compactifications --- arXiv:hep-th/0105097}

\textbf{Authors:} Steven B. Giddings, Shamit Kachru, Joseph Polchinski

\paragraph{主な主張}
Type IIBおよびF理論のコンパクト化で、量子化されたRR・NSフラックスにより大きなワープ階層を実現し、複素構造などのモジュライを固定できることを示した。 \cite{Giddings:2001yu}

\paragraph{新規性}
局所的なワープ喉部をコンパクトな弦背景へ組み込み、階層の大きさをフラックスの離散的選択と結び付けた。

\paragraph{位置付け}
フラックス安定化とワープ階層の基礎となるGKP構成であり、先頭次数で平坦なK\"ahler方向を残すno-scale構造を与える。

\paragraph{コメント（読書メモ）}
読書時の分類は「modular sym.なし」（この表現の適用範囲は未確認）。Type IIBのフラックスコンパクト化によるモジュライ安定化を扱う文献として読む。
